\subsection{Minimal projective resolutions}

Let M be an A-module and let

\[
\dots \longrightarrow \mathrm{P} _ {n} \xrightarrow {d _ {n}} \mathrm{P} _ {n - 1} \longrightarrow \dots \longrightarrow \mathrm{P} _ {0} \xrightarrow {d _ {0}} \mathrm{M} \longrightarrow 0\tag{P}
\]

a resolution of M. We say that (P) is a minimal projective resolution if, for every \(n \geqslant 0\) , the homomorphism \(\delta_{n}: P_{n} \to \operatorname{Im}(d_{n})\) induced by \(d_{n}\) is a projective cover (VIII, § 8, no. 5).

PROPOSITION 8. — Let M be an A-module, and let P and P' be two minimal projective resolutions of M, with f: P → P' a morphism of resolutions. Then f is an isomorphism. In particular, any two minimal projective resolutions of M are isomorphic.

Set \(\tilde{\mathbf{P}}_{n} = \mathbf{P}_{n}\) for \(n \neq -1\) and \(\tilde{\mathbf{P}}_{-1} = \mathbf{M}\) ; let us define similarly \(\tilde{\mathbf{P}}_{n}^{\prime}\) and set \(f_{-1} = 1_{\mathbf{M}}\) . Let us prove by induction starting from \(-1\) that \(f_{n}: \tilde{\mathbf{P}}_{n} \to \tilde{\mathbf{P}}_{n}^{\prime}\) is an isomorphism for all \(n\) . This is evident for \(n = -1\) ; suppose that \(f_{n}\) and \(f_{n-1}\) are isomorphisms. It follows from the commutativity of the diagram:

\[
\begin{array}{c c c} \mathbf {P} _ {n} & \xrightarrow {d _ {n}} & \mathbf {P} _ {n - 1} \\ \Big \downarrow f _ {n} & & \Big \downarrow f _ {n - 1} \\ \mathbf {P} _ {n} ^ {\prime} & \xrightarrow {d _ {n} ^ {\prime}} & \mathbf {P} _ {n - 1} ^ {\prime} \end{array}
\]

that \(f_{n}\) induces an isomorphism \(g_{n}\) of \(\operatorname{Ker} d_{n}\) onto \(\operatorname{Ker} d_{n}^{\prime}\) . It then follows from the commutativity of the diagram:

\[
\begin{array}{c c c} \mathrm{P} _ {n + 1} & \xrightarrow {\delta_ {n + 1}} & \text {Ker} d _ {n} \\ \Big \downarrow f _ {n + 1} & & \Big \downarrow g _ {n} \\ \mathrm{P} _ {n + 1} ^ {\prime} & \xrightarrow {\delta_ {n + 1} ^ {\prime}} & \text {Ker} d _ {n} ^ {\prime} \end{array}
\]

and from VIII, loc. cit., that \(f_{n+1}\) is an isomorphism.

COROLLARY. --- Let M be an A-module, and P and P' two projective resolutions of M; suppose that P is minimal. Let f : P → P' and g : P' → P be two morphisms of resolutions. Then f is injective, g is surjective, and P' is the direct sum of the subcomplexes Im f and Ker g. Moreover Ker g has zero homology.

Indeed, \(\alpha = g \circ f\) is an automorphism of P (prop. 8). Set \(\tilde{f} = f \circ \alpha^{-1}\) . We have

\[
\operatorname{Im} \tilde {f} = \operatorname{Im} f \quad \text { and } \quad g \circ \tilde {f} = 1 _ {\mathrm{P}},
\]

which shows that \(\mathbf{P}' = \operatorname{Im} \tilde{f} \oplus \operatorname{Ker} g\) . Since the sequence

\[
0 \rightarrow \operatorname{Ker} g \rightarrow \mathrm{P} ^ {\prime} \stackrel {g} {\rightarrow} \mathrm{P} \rightarrow 0
\]

is exact and \(g\) is a homologism, \(\operatorname{Ker} g\) has zero homology.

PROPOSITION 9. --- Let M be an A-module and (P, p) a projective resolution of M. Let r denote the radical of A. Suppose either that \(P_{n}\) is a finitely generated A-module for all n, or that r is nilpotent. Then for (P, p) to be minimal it is necessary and sufficient that the complex \((\mathbf{A}/\mathbf{r}) \otimes_{\mathbf{A}} \mathbf{P}\) have zero differential, in other words that

\[
d _ {n + 1} (\mathbf {P} _ {n + 1}) \subset \mathrm{r} \mathbf {P} _ {n} \quad \text { for all } n \geqslant 0.
\]

Suppose that (P, p) is minimal. By VIII, loc. cit., the homomorphism

\[
1 \otimes \delta_ {n}: (\mathrm{A} / \mathrm{r}) \otimes_ {\mathrm{A}} \mathrm{P} _ {n} \rightarrow (\mathrm{A} / \mathrm{r}) \otimes_ {\mathrm{A}} \operatorname{Im} d _ {n}
\]

is an isomorphism. It then follows from the exact sequence

\[
0 \longrightarrow \operatorname{Im} d _ {n + 1} \xrightarrow {j _ {n}} \mathrm{P} _ {n} \xrightarrow {\delta_ {n}} \operatorname{Im} d _ {n} \longrightarrow 0
\]

that the homomorphism \(1 \otimes j_n : (\mathbf{A} / \mathfrak{r}) \otimes_{\mathbf{A}} \operatorname{Im} d_{n+1} \to (\mathbf{A} / \mathfrak{r}) \otimes_{\mathbf{A}} \mathrm{P}_n\) is zero; since \(d_{n+1} = j_n \circ \delta_{n+1}\) , we deduce that \(1_{\mathbf{A} / \mathfrak{r}} \otimes d_{n+1} = 0\) for \(n \geqslant 0\) .

Conversely, suppose that for all \(n \geqslant 1\) , \(1 \otimes d_n\) is zero, in other words that \(\operatorname{Im} d_n = \operatorname{Ker} d_{n-1}\) is contained in \(\mathbf{rP}_{n-1}\) . Since \(\delta_{n-1}\) is surjective, it follows from VIII, loc. cit. that \(\delta_{n-1}\) is a projective cover for \(n \geqslant 1\) , hence that \((\mathbf{P}, p)\) is minimal.

PROPOSITION 10. --- Suppose that A is a left Noetherian local ring, and let M be a finitely generated A-module. Then M possesses a minimal resolution (P, p); for every \(n \geqslant 0\) , \(\mathbf{P}_{n}\) is a finitely generated free module.

Indeed, in the construction made in n° 5 (p.~53), one may, by virtue of VIII, loc. cit., take for \((\mathrm{L}_{0}, p)\) a projective cover of M, and for \(L_{n+1}\) a projective cover of Ker \(d_{n}\) . The resolution thus obtained is then minimal.

Remarks. --- 1) Let m denote the maximal ideal of A and set \(k = \mathbf{A}/\mathfrak{m}\) . Let P be a minimal projective resolution of M, and set \(b_{n} = \dim_{k}(k \otimes_{\mathbf{A}} \mathbf{P}_{n})\) . Then \(\mathbf{P}_{n}\) is a free A-module of rank \(b_{n}\) . It follows from the corollary of prop. 8 that for any other projective resolution \(\mathbf{P}'\) of M, one has \(\dim_{k}(k \otimes_{\mathbf{A}} \mathbf{P}_{n}') \geqslant b_{n}\) , and that equality holds if and only if \(\mathbf{P}'\) is minimal.

\begin{enumerate}
\def\labelenumi{\arabic{enumi})}
\setcounter{enumi}{1}
\tightlist
\item
  By prop. 9, \(b_n\) is the dimension over \(k\) of \(H_n(k \otimes_A P)\) , * in other words of \(\text{Tor}_n^A(k, M)\) . It is also the dimension over \(k\) of \(\text{Ext}_A^n(M, k)\) (cf.~X, p.~103, remark 3) *.
\end{enumerate}

